\section{Background}
\label{sec:background}

Argo floats park and drift at 1000m depth between profiling cycles.
Their locations are determined by currents at that depth, not by experimental design.
If 1000m currents are correlated with upper-ocean temperature or salinity, then Argo-based estimates of quantities like space-time averaged ocean heat content (OHC) could be biased.
This is widely suspected to be negligible, but the magnitude of this bias has not yet been bounded statistically.

Sampling of this kind---locations generated by a process correlated with the
field under study---is \emph{preferential} in the sense of \citet{diggle2010}:
the location process and the field are statistically dependent, so averages
formed from the observations need not match averages over the domain. Its effect
on areal averages and spatial prediction is documented across environmental
monitoring and ecology \citep{gelfand2012,zidek2014,conn2017,dumelle2022}, and
the standard response is a \emph{model-based correction}---a joint model for the
locations and the latent field, a shared Gaussian process in the canonical
treatment of \citet{diggle2010}, whose fit removes the bias. Such corrections
require a model for the sampling mechanism and a covariance estimated from the
very data whose locations are preferential, with the model risk both entail.

We take the complementary, design-based route: rather than correct the bias, we
ask first whether it needs correcting---how large the OHC preferential-sampling
bias can be, in the worst case, if one does nothing. The do-nothing estimator is
the naive area average, which inherits the entire discrepancy between where the
floats sample and the domain being averaged. We bound its bias, distribution-free
and then within a smoothness class for the OHC field, from the observed float
positions alone, the unknown field entering the guarantee only through a single,
a-priori-capped amplitude.

This worst case bounds the estimators actually used. OHC products interpolate the
profiles---by kriging, objective analysis, or optimal interpolation---and
integrate the interpolant \citep{baugh2022}; each is a linear functional of the
observations, of which the naive average is a suboptimal instance and kriging the
optimal one, so the worst-case bound \emph{provably upper-bounds the
preferential-sampling bias of any such interpolation estimator}
(Proposition~\ref{prop:rkhs-estimator}). It does not bound nonlinear, model-based
de-biasing, which uses the locations themselves as information and lies outside
this family; instead it makes such a correction's necessity assessable. If the
worst-case bias already sits within the interpolation error the products carry,
no correction---kriging-based or model-based---is warranted for OHC and its
attendant modelling can be foregone; if it does not, the bound quantifies exactly
what is at stake. The bound thus triages the design's susceptibility to
preferential sampling before any model is invoked to remove it.

This section sets up the estimand, the sampling mechanism, and the resulting sampling-induced bias.

\section{Methodology}
\label{sec:methodology}

All proofs are collected in Appendix~\ref{app:proofs}.

\subsection{Notation and Estimand}
\label{sec:def}

\begin{definition}[Space--time domain]
\label{def:domain}
Fix a spatial region of interest $A \subset S^{2}$, taken to be a Borel subset of the sphere with finite, positive surface area, and a study period $\T \subset \R$, a bounded interval.
The \emph{space--time domain} is the product set
$\X := A \times \T$, equipped with the product $\sigma$-algebra
$\F := \B(A)\otimes\B(\T)$, where $\B(A)$ and $\B(\T)$ are the Borel
$\sigma$-algebras on $A$ and $\T$. A generic element of $\X$ is written
$x = (s,t)$, with spatial location $s \in A$ and time $t \in \T$.
\end{definition}

\begin{definition}[Reference space--time measure]
\label{def:nu}
Let $\alpha$ be the normalized area measure on $A$, $\alpha(E)=|E|/|A|$,
and let $\tau$ be a probability measure on $\T$ specifying how the study
period is weighted. The \emph{reference} or \emph{target} probability
measure on $(\X,\F)$ is the product
\[
  \nu \;:=\; \alpha\otimes\tau .
\]
Unless stated otherwise $\tau$ is a seasonally balanced time-weighting.
When the window spans a whole number
of seasonal periods this is the uniform measure $\dd\tau=\dd t/|\T|$, and
$\nu$ reduces to normalized space--time volume, $\nu(E)=|E|/|\X|$.
\end{definition}

\begin{definition}[Sampling probability measure]
\label{def:pi}
Let $\pi$ be the \emph{sampling} probability measure on $(\X,\F)$.
\end{definition}


\begin{definition}[Target-population random profile]
\label{def:X}
Let $X = (S,T)$ be the \emph{coordinate map} on $(\X,\F,\nu)$, i.e., the identity function $X : \X \to \X$, $X(x) = x$.
Its distribution is the reference measure by construction, $X \sim \nu$.
Thus, $\Prob(X \in E) = \nu(E)$ for every $E \in \F$. 
\end{definition}

\begin{definition}[Sampled-population random profile]
\label{def:Z}
Let $Z = (S,T)$ be a random variable on $(\X,\F,\pi)$, with distribution the sampling measure, $Z \sim \pi$.
Thus, $\Prob(Z \in E) = \pi(E)$ for every $E \in \F$.
\end{definition}

$\nu(E)$ is the probability that a profile drawn \emph{uniformly} from the space-time domain lies in $E$.
Whereas, $\pi(E)$ is the probability that a profile drawn at random from the \emph{sampling} or realized advection process lies in $E$.

\begin{definition}[Observed profiles]
\label{def:Zi}
The $N$ observed profile locations are realizations of $N$ random elements $Z_1,\dots,Z_N$, each marginally distributed as $Z_i \sim \pi$ (Definition~\ref{def:Z}).
\end{definition}

\begin{definition}[Sampling weight]
\label{def:w}
Assume $\pi \ll \nu$. The \emph{sampling weight} is the Radon--Nikodym
derivative $w := \dd\pi/\dd\nu$: the $\nu$-almost-everywhere unique,
non-negative, $\F$-measurable function on $\X$ with
$\pi(E)=\int_E \dd\pi = \int_E w\dd\nu$ for all $E\in\F$. We assume further that
$w \in L^2(\nu)$.
\end{definition}

Composing with the target draw $X\sim\nu$ (Definition~\ref{def:X}) gives
the non-negative random variable $w(X)$; note that $w:\X\to[0,\infty)$ is
a deterministic function on the domain, while $w(X)$ is its composition
with the random location. Lebesgue differentiation gives
\[
  w(x) \;=\; \lim_{E \downarrow x} \frac{\pi(E)}{\nu(E)}
        \;=\; \lim_{E \downarrow x}
              \frac{\text{sampling mass of } E}{\text{reference mass of } E},
\]
so $w$ is the local ratio of sampling density to reference density. Where
$w \equiv 1$ the region is sampled exactly in proportion to its
space--time extent, i.e., as if uniformly; $w > 1$ marks over-sampling
(preferential), $w < 1$ under-sampling; and if $w = 0$ $\nu$-a.e. on a set
$E$ with $\nu(E)>0$, then $\pi(E)=0$ and $E$ is never sampled.

Intuitively, $w(X)$ reweights a target draw $X\sim\nu$ into a sampled draw $Z\sim\pi$.
Non-preferential sampling is $w\equiv 1$ ($\pi=\nu$).
\[
  \Prob(w(X) \in E) = \int_E w\dd\nu = \int_E \dd\pi = \pi(E).
\]


Because $\pi$ is a probability measure, $w(X)$ is mean-one,
\[
  \E_{\nu}\big[w(X)\big] \;=\; \int_{\X} w \dd\nu \;=\; \pi(\X) \;=\; 1,
\]
so $w-1$ is centered, and its dispersion under the reference measure,
\begin{equation}
  \sigma_w^2 \;:=\; \Var_{\nu}\big(w(X)\big)
  \;=\; \int_{\X} (w-1)^2 \dd\nu
  \;=\; \|w-1\|_{L^2(\nu)}^2
  \;=\; \int_{\X} w^2 \dd\nu - 1
  \;=\; \chi^2(\pi \,\|\, \nu),
  \label{eq:sigma-w}
\end{equation}
summarizes how inhomogeneous the sampling is: $\sigma_w = 0$ exactly when
$w = 1$ $\nu$-a.e., i.e. $\pi = \nu$ (non-preferential sampling), and
$\sigma_w$ grows as $\pi$ departs from $\nu$. Since $\E_\nu[w]=1$, the
scalar $\sigma_w$ is equivalently the coefficient of variation of $w(X)$,
a scale-free measure of sampling unevenness. Note that $w\in L^1(\nu)$ is
automatic, whereas the assumption $w\in L^2(\nu)$---equivalently
$\sigma_w<\infty$---is a genuine restriction, excluding sampling measures
so concentrated relative to $\nu$ that their $\chi^2$ divergence diverges.

Two features of \eqref{eq:sigma-w} shape what follows. It is defined
\emph{against the reference measure} $\nu$: it asks how far the sampling
law departs from the target-averaging weights. 

\begin{definition}[OHC random field]
\label{def:Y}
Let $Y: \Omega_Y \to L^2(\nu)$ be a random variable from ocean states to functional space.
\end{definition}

The bias we study is a property of the \emph{realized} ocean, so we condition on a realization of the field throughout.
Fix $\omega_Y\in\Omega_Y$ and write $y := Y(\omega_Y,\cdot)\in L^2(\nu)$ for the corresponding realization.
Because the study concerns a single physical ocean over the period $\T$, we identify $y$ with the true realized OHC field.
Conditional on the event $\{Y = y\}$; $y$ is a fixed element of $L^2(\nu)$.
Composing $y$ with the coordinate map yields the real-valued random variable $y(X): \X \to \R$ on the population space.

\begin{definition}[Estimand]
\label{def:mu}
The estimand is the \emph{realized mean OHC}, 
\[
  \mu = \E_\nu\big[y(X)\big] = \int_\X y \dd\nu.
\]
As a functional of the random field, $\mu(Y)$ is itself a random conditional on $\{Y=y\}$ it is a fixed scalar. 
Where
\[
  \sigma^2_Y(y) = \Var_{\nu}(y(X)) = \int_{\X}(y - \mu)^2\dd\nu = \|\,y - \mu\,\|^2_{L^2(\nu)}
\]
is the space--time variance of the realized field.
\end{definition}

The variability of $y$ over $\X$ is dominated by the seasonal cycle,
which inflates $\sigma_Y$---and hence every field factor in the bounds
below---while contributing nothing to $\mu$ itself over a window spanning
whole years. We therefore carry the field in decomposed form from the
outset.

\begin{definition}[Admissible offset and residual field]
\label{def:offset}
A function $m\in L^2(\nu)$ is an \emph{admissible offset} if it is fixed
(non-random), evaluable at every observed profile location, and its
reference integral $\E_\nu[m]$ is available in closed form. Given such an
$m$ we write
\[
  y' := y-m,\qquad
  \sigma_{Y'}^2 := \Var_\nu\big(y'(X)\big),\qquad
  D := \E_\pi[m]-\E_\nu[m]
\]
for the \emph{residual}, or \emph{deseasonalized}, field, its space--time
variance, and the \emph{offset aliasing}. The offset used throughout is
a harmonic model whose seasonal coefficients are fixed externally and
whose reference integral is exact.
\end{definition}

\begin{remark}[The offset does not change the estimand]
\label{rem:estimand-invariance}
Definition~\ref{def:mu} is stated for $y$ and is not altered by the
choice of $m$: since $\mu=\E_\nu[y']+\E_\nu[m]$ with the second term
known exactly, the estimand remains the realized mean OHC. We do
\emph{not} take $\E_\nu[y']$ as the target, which would make the
scientific quantity depend on an arbitrary modeling choice. What
deseasonalization changes is the \emph{estimator}
and, decisively, the field whose dispersion
enters every bound below: the operative factor is $\sigma_{Y'}$ rather
than $\sigma_Y$, with $\sigma_{Y'}=\sigma_Y\sqrt{1-\rho^2}$ for the
fraction $\rho^2$ of field variance the offset explains.
Over the analysis window, which spans an
integer number of years, the seasonal harmonics contribute exactly zero
to $\E_\nu[m]$, so nothing is
approximated here.
\end{remark}

\begin{definition}[Design mean]
\label{def:mu-w}
The \emph{design mean} is the mean of the realized field under the sampling measure,
\[
  \mu_w = \E_\pi\big[y(Z)\big] = \int_\X y \dd\pi.
\]
Equivalenty,
\[
  \mu_w = \int_\X w\,y \dd\nu = \E_\nu\big[w(X)\,y(X)\big].
\]

\end{definition}

Let $\hat\mu$ denote any estimator of $\mu$ built from the observed profiles (for example a kriging/Gaussian-process interpolant).
Its bias decomposes into a sampling term and a reconstruction term.

\subsection{Bias Decomposition}
\label{sec:defbias}

\begin{definition}[Bias decomposition]
\label{def:bias}
The bias of $\hat\mu$ decomposes as
\[
  B(y) := \hat\mu - \mu
  = \underbrace{\big(\mu_w - \mu\big)}_{\textstyle B_{\mathrm{samp}}(y)}
  + \underbrace{\big(\hat\mu - \mu_w\big)}_{\textstyle B_{\mathrm{pred}}(y)},
\]
where $B_{\mathrm{samp}}(y)$ is the \emph{preferential-sampling bias}
and $B_{\mathrm{pred}}(y)$ is the \emph{prediction bias}. \end{definition}

\begin{remark}[Scope of the guarantee]
\label{rem:bias-scope}
Definition~\ref{def:bias} is an algebraic identity, obtained by adding
and subtracting $\mu_w$, and so holds for every estimator. The bounds of
this paper apply to $B_{\mathrm{samp}}$ alone, which is a functional of
the design and the field and contains no estimator; everything specific
to $\hat\mu$ is absorbed into $B_{\mathrm{pred}}$. Consequently they
bound the \emph{total} bias of $\hat\mu$ precisely when
$\E[\hat\mu]=\mu_w$, so that $B_{\mathrm{pred}}$ is centred at zero. This
holds for the equal-weight sample mean over the observed profiles, and
for the deseasonalized control-variate estimator. For an estimator that reweights the
profiles---a kriging or objective-analysis interpolant, say---
$B_{\mathrm{pred}}$ is not centred and is not bounded here; the guarantee
then covers the sampling component only, and the reconstruction error
remains the responsibility of the interpolation method.
\end{remark}

The sampling term admits a covariance representation on the population
space, the starting point for the bound.

\begin{proposition}[Sampling-bias covariance identity]
\label{prop:cov-identity}
Conditional on $\{Y=y\}$ and the design, 
\[
  B_{\mathrm{samp}}(y) = \mu_w - \mu
  = \E_\nu\big[(w(X)-1)\,y(X)\big]
  = \Cov_\nu\big(w(X),\,y(X)\big).
\]
\end{proposition}
\deferproof{cov-identity}{%
\begin{proof}[Proof of Proposition~\ref{prop:cov-identity}]
By Definitions~\ref{def:mu} and~\ref{def:mu-w},
$\mu_w-\mu=\E_\nu[w(X)y(X)]-\E_\nu[y(X)]=\E_\nu[(w(X)-1)y(X)]$. Expanding the
covariance about $\mu$ and using~$\E_{\nu}[w(X)]=1$,
\begin{align*}
  \Cov_\nu(w(X),y(X)) &= \E_\nu\big[(w(X)-1)(y(X)-\mu)\big] \\
  &= \E_\nu\big[(w(X)-1)y(X)\big] - \mu\,\E_\nu[w(X)-1] \\
  &= \E_\nu\big[(w(X)-1)y(X)\big],
\end{align*}
since $\E_\nu[w(X)-1]=0$.
\end{proof}%
}

\begin{remark}[No sampling bias under uncorrelated design]
\label{rem:zero-bias}
If, for the realized field, the weight and the field are uncorrelated
under $\nu$---in particular whenever $w\equiv1$---then
$\Cov_\nu(w,y)=0$ and $B_{\mathrm{samp}}(y)=0$. Preferential-sampling
bias arises precisely from co-variation between where the floats
concentrate ($w$) and the value of the field ($y$).
\end{remark}

\begin{proposition}[Distribution-free $L^2$ bound]
\label{prop:l2-bound}
Under the conditions of Proposition~\ref{prop:cov-identity},
\[
  \big|B_{\mathrm{samp}}(y)\big|
  = \big|\Cov_\nu(w(X),y(X))\big|
  = |\rho_{w,y}|\,\sigma_w\,\sigma_y
  \;\le\; \sigma_w\,\sigma_y,
\]
where $\rho_{w,y}\in[-1,1]$ is the $\nu$-correlation between $w(X)$ and
$y(X)$. The bound requires only finite second moments and imposes no
smoothness assumption on the field.
\end{proposition}

\begin{remark}[Relation to preferential-sampling and survey-sampling inference]
\label{rem:related-samp}
The identity $B_{\mathrm{samp}}=\mu_w-\mu=\Cov_\nu(w,y)$ is the design-based
signature of preferential sampling. The design mean $\mu_w$ is the expectation
of the naive area average $N^{-1}\sum_i y(Z_i)$ under the sampling law, so
$B_{\mathrm{samp}}$ is that estimator's bias---non-zero exactly when the
location process and the field are dependent, which is the defining condition of
preferential sampling \citep{diggle2010}. The same covariance is the bias of the
unweighted mean under unequal inclusion probabilities in design-based survey
sampling, removed there by inverse-probability (Horvitz--Thompson) weighting
\citep{horvitz1952,sarndal1992}; its effect on spatial prediction and on areal
averages is documented across environmental monitoring and ecology
\citep{gelfand2012,zidek2014,conn2017,dumelle2022}. That literature is
predominantly \emph{model-based}---it posits a joint model for the locations and
the latent field and corrects the bias by fitting it, as in the shared-latent
Gaussian-process remedy of \citet{diggle2010}. We instead treat the same
functional $\Cov_\nu(w,y)$ \emph{design-based} and \emph{bound} it:
distribution-free in Proposition~\ref{prop:l2-bound}, and within a smoothness
class in \S\ref{sec:rkhs}, so that no model for the sampling mechanism or the
field enters the guarantee.
\end{remark}

\section{Estimation}
\label{sec:estimation}

\begin{definition}[Equal-measure partition and the projection operators]
\label{def:partition}
A \emph{level-$K$ partition} is a collection
$\mathcal{P}_K=\{C_1,\dots,C_K\}$ of disjoint measurable cells with
$\bigcup_j C_j=\X$ and equal reference measure $\nu(C_j)=1/K$ for all
$j$. Write $\PK:L^2(\nu)\to L^2(\nu)$ for the \emph{cell-averaging
operator}, i.e., the conditional expectation
$\PK f=\E_\nu[f\mid\mathcal{P}_K]$, which replaces a function by its cell
averages,
\[
  (\PK f)(x)=\bar f_j:=\frac{1}{\nu(C_j)}\int_{C_j} f\dd\nu
  \qquad\text{for }x\in C_j ,
\]
and write $\PKp:=\mathsf{I}-\PK$ for the complementary
\emph{within-cell} (or \emph{fluctuation}) operator,
\[
  (\PKp f)(x) := f(x)-(\PK f)(x) = f(x)-\bar f_j
  \qquad\text{for }x\in C_j .
\]
We abbreviate $f^{\perp}:=\PKp f$, so that every $f\in L^2(\nu)$ splits
as
\begin{equation}
  f=\PK f+\PKp f,\qquad\text{i.e.}\qquad \mathsf{I}=\PK+\PKp ,
  \label{eq:split}
\end{equation}
into its between-cell and within-cell parts.
\end{definition}

The two operators project onto complementary subspaces of $L^2(\nu)$.
Let
\[
  V_K=\Big\{\textstyle\sum_{j=1}^{K}a_j\mathbf{1}_{C_j}\;:\;a_j\in\R\Big\}
\]
be the $K$-dimensional subspace of functions constant on each cell (the
\emph{resolution-$K$} functions), whose indicator basis
$\{\mathbf{1}_{C_1},\dots,\mathbf{1}_{C_K}\}$ is orthogonal because the
cells are disjoint. Its orthogonal complement is
\[
  V_K^{\perp}=\Big\{g\in L^2(\nu)\;:\;\int_{C_j}g\dd\nu=0
    \ \text{for every }j\Big\},
\]
the functions with zero average on every cell, which carry only sub-cell
detail.

\begin{remark}
\label{rem:buildK}
Assuming $\nu$ is a measure over normalized area and time, equal-measure
cells are constructed by partitioning  
longitude into equal intervals, $\sin(\mathrm{lat})$ into equal
intervals (equal area on the sphere), and time into equal intervals,
so that $C_j$ is a longitude~$\times\sin(\mathrm{lat})\times$ time box
and $K=K_{\mathrm{lon}}K_{\mathrm{lat}}K_{\mathrm{t}}$. Each cell then
has equal meausre $\nu(C_j)=1/K$.
\end{remark}

\begin{lemma}[$\PK$ and $\PKp$ are complementary orthogonal projections]
\label{lem:projection}
For every $f\in L^2(\nu)$ we have $\PK f\in V_K$ and $\PKp f\in
V_K^{\perp}$; the operator $\PK$ is the orthogonal projection of
$L^2(\nu)$ onto $V_K$, and $\PKp$ the orthogonal projection onto
$V_K^{\perp}$. Consequently both are idempotent and self-adjoint,
\[
  \PK^2=\PK,\qquad (\PKp)^2=\PKp,\qquad \PK\PKp=\PKp\PK=0,
\]
their ranges are orthogonal, $\langle \PK f,\ \PKp g\rangle_\nu=0$ for all
$f,g\in L^2(\nu)$, and the Pythagorean identity
\begin{equation}
  \|f\|_{L^2(\nu)}^2=\|\PK f\|_{L^2(\nu)}^2+\|\PKp f\|_{L^2(\nu)}^2
  \label{eq:pythagoras}
\end{equation}
holds.
\end{lemma}
\deferproof{projection}{%
\begin{proof}[Proof of Lemma~\ref{lem:projection}]
$\PK f$ is constant on each cell, so $\PK f\in V_K$. For the residual,
use $\bar f_j\,\nu(C_j)=\int_{C_j}f\dd\nu$ from the definition of the
cell mean:
\[
  \int_{C_j}\PKp f\dd\nu
  =\int_{C_j}\big(f-\bar f_j\big)\dd\nu
  =\int_{C_j}f\dd\nu-\bar f_j\,\nu(C_j)=0
  \qquad\text{for every }j,
\]
so $\PKp f\in V_K^{\perp}$ by the characterization above. Orthogonality
of the ranges follows cell by cell: $\PK f$ takes the constant value
$\bar f_j$ on $C_j$, so
\[
  \big\langle \PK f,\ \PKp g\big\rangle_\nu
  =\sum_{j=1}^{K}\int_{C_j}\bar f_j\,(\PKp g)\dd\nu
  =\sum_{j=1}^{K}\bar f_j\underbrace{\int_{C_j}\PKp g\dd\nu}_{=\,0}
  =0 .
\]
In particular $f-\PK f$ is orthogonal to all of $V_K$, which is the
defining property of the orthogonal projection onto $V_K$; since
$\PKp=\mathsf{I}-\PK$, it is the orthogonal projection onto
$V_K^{\perp}$. Idempotence, self-adjointness and $\PK\PKp=0$ are the
standard consequences, and applying the orthogonality to
$f=\PK f+\PKp f$ gives \eqref{eq:pythagoras}.
\end{proof}%
}

Intuitively, subtracting the cell mean removes everything a
piecewise-constant function could capture: the average has already
extracted all of $f$ that is visible at cell resolution, so the residual
has nothing left to say there. That residual is $\PKp f$, the
\emph{within-cell fluctuation}: on each cell it equals $f$ minus that
cell's mean, and hence integrates to zero within every cell. It lives in
$V_K^{\perp}$, the functions carrying only sub-cell detail, so the split
$f=\PK f+\PKp f$ of~\eqref{eq:split} separates $f$ into its coarse,
between-cell part and its fine, within-cell part.

\subsection{Decomposing the weight and the field}
\label{sec:decomp-wy}

Lemma~\ref{lem:projection} applies to any $f\in L^2(\nu)$; the two
functions the bound requires are the centered weight $w-1$ and the field
anomaly $y-\mu$. Both are centered under $\nu$: $\E_\nu[w(X)]=1$ by
\eqref{eq:sigma-w}, and $\E_\nu[y(X)]=\mu$ by Definition~\ref{def:mu}.
Applying the operator identity~\eqref{eq:split} to each,
\begin{align}
  w-1 &= \underbrace{\PK(w-1)}_{\text{between-cell}}
        \;+\;\underbrace{\PKp(w-1)}_{\text{within-cell}}
        \;=\;w^{\star}+w^{\perp},
  \label{eq:w-split}\\[2pt]
  y-\mu &= \underbrace{\PK(y-\mu)}_{\text{between-cell}}
        \;+\;\underbrace{\PKp(y-\mu)}_{\text{within-cell}}
        \;=\;y^{\star}+y^{\perp},
  \label{eq:y-split}
\end{align}
where we have written
\[
  w^{\star}:=\PK(w-1)=\PK w-1,\qquad w^{\perp}:=\PKp(w-1)=\PKp w,
\]
\[
  y^{\star}:=\PK(y-\mu)=\PK y-\mu,\qquad y^{\perp}:=\PKp(y-\mu)=\PKp y,
\]
using that $\PK$ fixes constants ($\PK 1=1$, since $1\in V_K$) and hence
$\PKp$ annihilates them ($\PKp 1=0$). Concretely, the resolved parts
$w^{\star},y^{\star}$ are piecewise constant, equal to $\bar w_j-1$ and
$\bar y_j-\mu$ on $C_j$, while the unresolved parts $w^{\perp},y^{\perp}$
are $w-\bar w_j$ and $y-\bar y_j$ on $C_j$, each with zero average on
every cell.

\subsubsection*{Decomposition of the variances}

Because the two components in \eqref{eq:w-split} are orthogonal
(Lemma~\ref{lem:projection}), the Pythagorean
identity~\eqref{eq:pythagoras} splits the dispersion of the weight into
between- and within-cell pieces. Define
\begin{equation}
  \sigma_{w^{\star}} := \big\|w^{\star}\big\|_{L^2(\nu)},
  \qquad
  \sigma_{w^{\perp}} := \big\|w^{\perp}\big\|_{L^2(\nu)},
  \label{eq:sigma-w-parts}
\end{equation}
and analogously for the field,
\begin{equation}
  \sigma_{Y^{\star}} := \big\|y^{\star}\big\|_{L^2(\nu)},
  \qquad
  \sigma_{Y^{\perp}} := \big\|y^{\perp}\big\|_{L^2(\nu)} .
  \label{eq:sigma-Y-parts}
\end{equation}
Then
\begin{equation}
  \sigma_w^2=\sigma_{w^{\star}}^2+\sigma_{w^{\perp}}^2,
  \qquad
  \sigma_Y^2=\sigma_{Y^{\star}}^2+\sigma_{Y^{\perp}}^2 .
  \label{eq:var-split}
\end{equation}

These are exactly the law of total variance. Since
$\PK f=\E_\nu[f\mid\mathcal{P}_K]$, writing $X\sim\nu$ and letting
$J=J(X)$ denote the (random) index of the cell containing $X$,
\eqref{eq:var-split} reads
\[
  \underbrace{\Var_\nu\big(w(X)\big)}_{\sigma_w^2}
  =\underbrace{\Var_\nu\Big(\E_\nu\big[w(X)\mid J\big]\Big)}_{
      \sigma_{w^{\star}}^2\ =\ \frac1K \sum_j (\bar w_j-1)^2}
  \;+\;
  \underbrace{\E_\nu\Big[\Var_\nu\big(w(X)\mid J\big)\Big]}_{
      \sigma_{w^{\perp}}^2\ =\ \frac1K \sum_j \Var_\nu(w\mid C_j)} ,
\]
and identically for $y$: the resolved term is the variance of the cell
means, and the unresolved term is the mean of the within-cell variances.
In words, $\sigma_{w^{\star}}$ measures how unevenly sampling effort is
distributed \emph{across} cells---the dispersion of the cell weights
$\bar w_j=K\pi_j$ about one---while $\sigma_{w^{\perp}}$ measures the
sampling structure \emph{inside} cells, which the partition cannot see.
The same split separates field variation across cells from field
variation within them.

\subsubsection*{Why the split is useful}

The decompositions \eqref{eq:w-split}--\eqref{eq:y-split} are what allow
the bias itself to be separated by scale. Substituting them into the
covariance identity of Proposition~\ref{prop:cov-identity} and using the
orthogonality $\langle \PK f,\PKp g\rangle_\nu=0$ of
Lemma~\ref{lem:projection} to kill the two cross terms,
\begin{equation}
  B_{\mathrm{samp}}
  =\big\langle\, w-1,\; y-\mu \,\big\rangle_\nu
  =\underbrace{\big\langle\, w^{\star},\; y^{\star} \,\big\rangle_\nu}_{\text{between-cell}}
  \;+\;\underbrace{\big\langle\, w^{\perp},\; y^{\perp} \,\big\rangle_\nu}_{\text{within-cell}} ,
  \label{eq:bias-split}
\end{equation}
so the resolved and unresolved scales contribute additively and without
interference. Applying Cauchy--Schwarz to each term of
\eqref{eq:bias-split} separately pairs $\sigma_{w^{\star}}$ with
$\sigma_{Y^{\star}}$ and $\sigma_{w^{\perp}}$ with $\sigma_{Y^{\perp}}$,
giving
\[
  \big|B_{\mathrm{samp}}\big|
  \;\le\;\sigma_{w^{\star}}\,\sigma_{Y^{\star}}
  \;+\;\sigma_{w^{\perp}}\,\sigma_{Y^{\perp}} ,
\]
which is the content of Proposition~\ref{prop:scale-decomp}. This is what
makes a coarse partition---one the float counts can actually
support---admissible for bounding the bias, provided the within-cell
field variation $\sigma_{Y^{\perp}}$ is small.


\begin{proposition}[Scale-decomposition bound]
\label{prop:scale-decomp}
Conditional on the realized field $\{Y=y\}$ and the design, and under
the hypotheses of Proposition~\ref{prop:cov-identity}, for any
partition $\mathcal{P}_K$,
\[
  B_{\mathrm{samp}}
  =\underbrace{\big\langle\,w^{\star},\ y^{\star}\,\big\rangle_\nu}_{\text{between-cell}}
  +\underbrace{\big\langle\,w^{\perp},\ y^{\perp}\,\big\rangle_\nu}_{\text{within-cell}},
\]
and
\[
  |B_{\mathrm{samp}}|
  \;\le\; \sigma_{w^{\star}}\,\sigma_{Y^{\star}}
        \;+\;\sigma_{w^{\perp}}\,\sigma_{Y^{\perp}}
  \;\le\; \sigma_w\,\sigma_Y ,
\]
with $w^{\star},w^{\perp},y^{\star},y^{\perp}$ as in
\eqref{eq:w-split}--\eqref{eq:y-split} and the four dispersions as in
\eqref{eq:sigma-w-parts}--\eqref{eq:sigma-Y-parts}. In particular, the
scale-decomposed bound is never looser than the raw Cauchy--Schwarz
bound $\sigma_w\sigma_Y$.
\end{proposition}
\deferproof{scale-decomp}{%
\begin{proof}[Proof of Proposition~\ref{prop:scale-decomp}]
Substitute the decompositions $w-1=w^{\star}+w^{\perp}$ and
$y-\mu=y^{\star}+y^{\perp}$ of \eqref{eq:w-split}--\eqref{eq:y-split}
into the covariance identity
$B_{\mathrm{samp}}=\langle w-1,\,y-\mu\rangle_\nu$
(Proposition~\ref{prop:cov-identity}). Bilinearity of the inner product
gives four terms,
\begin{align*}
  B_{\mathrm{samp}}
  &=\big\langle\, w^{\star}+w^{\perp},\ y^{\star}+y^{\perp} \,\big\rangle_\nu\\[2pt]
  &=\underbrace{\big\langle\, w^{\star},\ y^{\star} \,\big\rangle_\nu}_{\text{between-cell}}
   \;+\;\underbrace{\big\langle\, w^{\star},\ y^{\perp} \,\big\rangle_\nu}_{\text{cross term}}
   \;+\;\underbrace{\big\langle\, w^{\perp},\ y^{\star} \,\big\rangle_\nu}_{\text{cross term}}
   \;+\;\underbrace{\big\langle\, w^{\perp},\ y^{\perp} \,\big\rangle_\nu}_{\text{within-cell}} .
\end{align*}
The two cross terms vanish by the orthogonality $V_K\perp V_K^{\perp}$ of
Lemma~\ref{lem:projection}: $w^{\star},y^{\star}\in V_K$ are constant on
each cell, while $w^{\perp},y^{\perp}\in V_K^{\perp}$ have zero average
on each cell. Explicitly, splitting the integral over the cells and
pulling out the constant value $\bar w_j-1$ of $w^{\star}$ on $C_j$,
\[
  \big\langle\, w^{\star},\ y^{\perp} \,\big\rangle_\nu
  =\int_{\X} w^{\star}\,y^{\perp}\dd\nu
  =\sum_{j=1}^{K}\int_{C_j}(\bar w_j-1)\,y^{\perp}\dd\nu
  =\sum_{j=1}^{K}(\bar w_j-1)\underbrace{\int_{C_j}y^{\perp}\dd\nu}_{=\,0}
  =0 ,
\]
and symmetrically $\langle w^{\perp},y^{\star}\rangle_\nu=0$, exchanging
the roles of $w$ and $y$. Only the between-cell and within-cell terms
survive, which is the stated two-term identity; Cauchy--Schwarz applied
to each of them separately yields
\[
  |B_{\mathrm{samp}}|
  \;\le\;\big|\big\langle w^{\star},y^{\star}\big\rangle_\nu\big|
        +\big|\big\langle w^{\perp},y^{\perp}\big\rangle_\nu\big|
  \;\le\;\sigma_{w^{\star}}\sigma_{Y^{\star}}
        +\sigma_{w^{\perp}}\sigma_{Y^{\perp}} .
\]

For the final inequality, apply the Cauchy--Schwarz inequality in $\R^2$
to the vectors
$\mathbf a=(\sigma_{w^{\star}},\sigma_{w^{\perp}})$ and
$\mathbf b=(\sigma_{Y^{\star}},\sigma_{Y^{\perp}})$:
\[
  \sigma_{w^{\star}}\sigma_{Y^{\star}}+\sigma_{w^{\perp}}\sigma_{Y^{\perp}}
  =\mathbf a\cdot\mathbf b
  \;\le\;\|\mathbf a\|\,\|\mathbf b\|
  =\sqrt{\sigma_{w^{\star}}^2+\sigma_{w^{\perp}}^2}\,
   \sqrt{\sigma_{Y^{\star}}^2+\sigma_{Y^{\perp}}^2}
  =\sigma_w\,\sigma_Y ,
\]
the last equality being the Pythagorean identities \eqref{eq:var-split}.
Equality holds precisely when $\mathbf a$ and $\mathbf b$ are parallel,
i.e. when $w$ and $y$ distribute their dispersion between the resolved
and unresolved scales in the same proportion,
$\sigma_{w^{\perp}}/\sigma_{w^{\star}}=\sigma_{Y^{\perp}}/\sigma_{Y^{\star}}$;
otherwise the scale-decomposed bound is strictly tighter.
\end{proof}%
}

Finally, we record how both decompositions behave under a change of
partition, since the resolution at which they are evaluated is itself a
choice.

\begin{lemma}[Coarsening is a monotone lower bound]
\label{lem:coarsening}
For any partition $\mathcal{P}_K$,
\[
  \sigma_{w^{\star}}\;\le\;\sigma_w,
  \qquad\text{with}\qquad
  \sigma_w^2-\sigma_{w^{\star}}^2
  =\sigma_{w^{\perp}}^2
  =\E_\nu\big[\Var_\nu(w\mid\mathcal{P}_K)\big]\;\ge\;0 .
\]
Moreover, under refinement the bound tightens monotonically: if
$\mathcal{P}_{K'}$ refines $\mathcal{P}_{K}$, in the sense that every
cell of $\mathcal{P}_{K}$ is a union of cells of $\mathcal{P}_{K'}$, then
$\sigma_{w^{\star}}(\mathcal{P}_{K})\le\sigma_{w^{\star}}(\mathcal{P}_{K'})\le\sigma_w$.
The same statements hold verbatim for the field, with
$(\sigma_{Y^{\star}},\sigma_{Y^{\perp}},\sigma_Y)$ in place of
$(\sigma_{w^{\star}},\sigma_{w^{\perp}},\sigma_w)$.
\end{lemma}

\deferproof{coarsening}{%
\begin{proof}
The identity is the Pythagorean decomposition \eqref{eq:var-split},
$\sigma_w^2=\sigma_{w^{\star}}^2+\sigma_{w^{\perp}}^2$, together with the
identification of the within-cell term as an expected conditional
variance (Section~\ref{sec:decomp-wy}); non-negativity of
$\sigma_{w^{\perp}}^2=\|w^{\perp}\|_{L^2(\nu)}^2$ gives
$\sigma_{w^{\star}}\le\sigma_w$. For monotonicity, refinement enlarges the
subspace of piecewise-constant functions, $V_{K}\subseteq V_{K'}$, and
the norm of an orthogonal projection cannot decrease when the range is
enlarged: $\|\mathsf{P}_{V_{K}}f\|\le\|\mathsf{P}_{V_{K'}}f\|$ for every
$f\in L^2(\nu)$. Applying this to $f=w-1$ gives the claim.
\end{proof}%
}

\section{Implementation}
\label{sec:implementation}

The variance components that define the preferential sampling bound can be estimated without fitting models.
The between-cell weight dispersion $\sigma_{w^{\star}}$ follows from
float counts alone (Section~\ref{sec:sigmawK}). The leading field factor
is bounded above by the observed OHC variance together with a
conservative kurtosis constant (Section~\ref{sec:sigmaY-bound}), applied
to the deseasonalized field $y'$---the seasonal cycle having been
removed as a control variate whose own sampling bias is computed rather
than bounded. The residual factor
$\sigma_{Y'^{\perp}}$ is estimated from within-cell float pairs. No external field product enters the
guarantee as an amplitude; where one is consulted it supplies a scale or
a dimensionless ratio only (Remark~\ref{rem:conservative-Y}).
$\sigma_{w^{\perp}}$, the sampling variability \emph{within} the cell scale cannot be estimated from a sparse point pattern without reintroducing a within-cell intensity model.
However, proposition~\ref{prop:scale-decomp} confines it to a product with $\sigma_{Y^{\perp}}$, so the magnitude of $\sigma_{w^{\star}}$ only matters to the extent the OHC field varies within cells.
Therefore, our goal is to determine the cell scale such that within-cell OHC variability, $\sigma_{Y^{\perp}}$, is small such that the $\sigma_{w^{\star}}$ $\sigma_{Y^{\perp}}$ product is negligible.

\subsection{Estimating $\sigma_{w^{\star}}$ from float counts}
\label{sec:sigmawK}

\begin{definition}[Cell sampling weight]
\label{def:sweight}
Let $\bar w_j$ be the \emph{sampling weight} for cell $j$, which is the ratio of the sampling mass of the cell to its reference mass,
\begin{equation}
  \bar w_j \;=\; \frac{1}{\nu(C_j)}\int_{C_j}w\dd\nu \;=\; \frac{\pi(C_j)}{\nu(C_j)} \;=\; \frac{\pi_j}{1/K}
  \;=\; K\pi_j ,
  \label{eq:cell-weight}
\end{equation}
where $\pi_j:=\pi(C_j)$ for the probability that a randomly sampled profile lies in cell $C_j$.
\end{definition}

This is coarse cell counterpart of the pointwise weight $w$ of
Definition~\ref{def:w}: $\bar w_j=1$ means cell $j$ is sampled exactly as
under uniform coverage, $\bar w_j>1$ over-sampled (preferential),
$\bar w_j<1$ under-sampled, and $\bar w_j=0$ unsampled. Because
$\sum_j\pi_j=1$ and $\nu(C_j)=1/K$, the weights are mean-one under the
reference measure, $\sum_j\nu(C_j)\bar w_j=\sum_j\pi_j=1$.

Let $n_j:=N(C_j)$ be the number of profiles observed in cell $j$ and
$N=\sum_j n_j$ the total. The natural estimator of $\bar w_j$ is
\begin{equation}
  \widehat{\bar w}_j=K\,\frac{n_j}{N}=\frac{n_j}{\,N/K\,}
  =\frac{\text{observed count in }C_j}{\text{count expected under uniform}} ,
  \label{eq:wbar-hat}
\end{equation}
so that a map of $\widehat{\bar w}_j$ over the domain displays directly
where floats over- and under-sample relative to uniform coverage.

\begin{lemma}[Cell sampling weight variance]
\label{lem:sigmawcell}
The variance of the cell sampling weights is
\[
  \sigma^2_{w^{\star}} = K\sum_{j=1}^{K}\pi_j^2-1 .
\]
\end{lemma}
\deferproof{sigmawcell}{%
\begin{proof}[Proof of Lemma~\ref{lem:sigmawcell}]
$\sigma^2_{w^{\star}}$ is the variance of the cells weights about one.
This is because $w^{\star}=\PK w-1$ is a constant function on each cell with cell value on $C_j$ is $\bar w_j-1$.

Accordingly,
\begin{align*}
  \sigma_{w^{\star}}^2
  &=\int_{\X}(w^{\star})^2\dd\nu
   =\sum_{j=1}^{K}\int_{C_j}\big(\bar w_j-1\big)^2\dd\nu
   =\sum_{j=1}^{K}\big(\bar w_j-1\big)^2\,\nu(C_j)
   =\frac{1}{K}\sum_{j=1}^{K}\big(K\pi_j-1\big)^2 \\
  &=\frac{1}{K}\sum_{j=1}^{K}\Big(K^2\pi_j^2-2K\pi_j+1\Big)
   =K\sum_{j=1}^{K}\pi_j^2\;-\;2\underbrace{\sum_{j=1}^{K}\pi_j}_{=\,1}
     \;+\;\frac{1}{K}\underbrace{\sum_{j=1}^{K}1}_{=\,K} \\
  &=K\sum_{j=1}^{K}\pi_j^2-1 .
\end{align*}
\end{proof}%
}

Note, uniform sampling within the partition ($\pi_j=1/K$ for all $j$, i.e.
$\bar w_j\equiv1$) gives $K\cdot K^{-1}-1=0$.
Whereas when all sampling mass is in a single cell gives $K-1$.
Hence, $0\le\sigma_{w^{\star}}\le\sqrt{K-1}$.
Thus, between-cell sampling variability can increase arbitrarily as the number of cells $K$ grows, degenerating as almost all cells have zero observed counts.

\begin{assumption}[Profiles sample the sampling measure]
\label{ass:sample-pi}
Conditional on the total $N$, the random elements $Z_1,\dots,Z_N$ of
Definition~\ref{def:Zi} are independent draws from the sampling measure
$\pi$ of Definition~\ref{def:pi}; equivalently, the cell counts follow
$(n_1,\dots,n_K)\sim\mathrm{Multinomial}\big(N;\pi_1,\dots,\pi_K\big)$
with $\pi_j=\pi(C_j)$.
\end{assumption}

The plug-in $\hat\pi_j=n_j/N$ gives
\[
  \sigwhatsqplug
  =K\sum_j\hat\pi_j^2-1=\frac{X^2}{N},
  \qquad X^2=\sum_{j=1}^{K}\frac{(n_j-N/K)^2}{N/K},
\]
where $X^2$ is \emph{Pearson's goodness-of-fit statistic for
uniformity}.
This plug-in estimator is upward biased and removing the diagonal (self-pairs) removes the bias (Lemma~\ref{lem:unbiased}):
\begin{equation}
  \sigwhatsq
  \;=\;K\sum_{j=1}^{K}\frac{n_j(n_j-1)}{N(N-1)}\;-\;1
  \label{eq:sigmawK-hat}
\end{equation}

\begin{remark}[Negative estimates]
\label{rem:negative}
Being unbiased for a non-negative quantity,
\eqref{eq:sigmawK-hat} can take negative values in finite samples.
A negative value arises exactly when the observed coincidence rate falls below the uniform rate, $\sum_j n_j(n_j-1)/[N(N-1)]<1/K$, and so indicates sampling at least as
even as the reference measure.
We retain the signed value for interval construction and report
$\sigwhat=\sqrt{\max\{0,\sigwhatsq\}}$
as the point estimate; near-zero results are stated on the $\sigma^2$
scale with their signed interval, as ``consistent with zero'' rather than
truncated.
\end{remark}

\begin{lemma}[Bias of the plug-in; unbiasedness of \eqref{eq:sigmawK-hat}]
\label{lem:unbiased}
Under Assumption~\ref{ass:sample-pi},
\[
  \E\big[\sigwhatsqplug\big]
  =\sigma_{w^{\star}}^2+\frac{K}{N}\Big(1-\textstyle\sum_j\pi_j^2\Big),
  \qquad
  \E\big[\sigwhatsq\big]=\sigma_{w^{\star}}^2 .
\]
The plug-in bias is non-negative, vanishing only in the degenerate case
$\sum_j\pi_j^2=1$ (all sampling mass in one cell), and grows with $K/N$;
the diagonal-removed estimator~\eqref{eq:sigmawK-hat} is exactly unbiased
for every $N\ge 2$ and every $K$.
\end{lemma}
\deferproof{unbiased}{%
\begin{proof}[Proof of Lemma~\ref{lem:unbiased}]
\begin{align}
  \mathrm{Bias}\big(\sigwhatsqplug\big)
  &=\E\Big[K\sum_{j=1}^{K}\hat\pi_j^2-1\Big]
    -\Big(K\sum_{j=1}^{K}\pi_j^2-1\Big)
   =K\sum_{j=1}^{K}\Big(\E\big[\hat\pi_j^2\big]-\pi_j^2\Big), \label{eq:bias-plug}\\
  \mathrm{Bias}\big(\sigwhatsq\big)
  &=\E\Big[K\sum_{j=1}^{K}\frac{n_j(n_j-1)}{N(N-1)}-1\Big]
    -\Big(K\sum_{j=1}^{K}\pi_j^2-1\Big)
   =K\sum_{j=1}^{K}\Big(\E\Big[\frac{n_j(n_j-1)}{N(N-1)}\Big]-\pi_j^2\Big).
  \label{eq:bias-diag}
\end{align}
It therefore suffices to evaluate the two cell-level expectations.

\emph{Step 1: moments of the counts.} Under
Assumption~\ref{ass:sample-pi} the counts are multinomial, so each
marginal is $n_j\sim\mathrm{Binomial}(N,\pi_j)$, giving
\[
  \E[n_j]=N\pi_j,\qquad \Var(n_j)=N\pi_j(1-\pi_j),\qquad
  \E[n_j^2]=\Var(n_j)+\E[n_j]^2=N\pi_j(1-\pi_j)+N^2\pi_j^2 .
\]

\emph{Step 2: the plug-in cell term is biased.} Dividing the second
moment by $N^2$,
\[
  \E\big[\hat\pi_j^2\big]=\frac{\E[n_j^2]}{N^2}
  =\pi_j^2+\frac{1}{N}\,\pi_j(1-\pi_j)
  \;>\;\pi_j^2 \quad\text{whenever }0<\pi_j<1 ,
\]
so each cell term over-estimates $\pi_j^2$ by $\pi_j(1-\pi_j)/N$.
Substituting into \eqref{eq:bias-plug} and using $\sum_j\pi_j=1$,
\[
  \mathrm{Bias}\big(\sigwhatsqplug\big)
  =K\sum_{j=1}^{K}\frac{\pi_j(1-\pi_j)}{N}
  =\frac{K}{N}\Big(\sum_{j=1}^{K}\pi_j-\sum_{j=1}^{K}\pi_j^2\Big)
  =\frac{K}{N}\Big(1-\sum_{j=1}^{K}\pi_j^2\Big),
\]
which is the first claim. It is non-negative because
$\sum_j\pi_j^2\le\big(\sum_j\pi_j\big)^2=1$, vanishing only when all
sampling mass lies in a single cell, and it grows linearly in $K/N$.

\emph{Step 3: removing the self-pairs.} The inflation arises because
$n_j^2$ counts all $N^2$ \emph{ordered} pairs of profiles falling in
$C_j$, including the $N$ pairs of a profile with itself. A pair of
\emph{distinct} profiles lies in $C_j$ with probability $\pi_j^2$, since
the two are independent draws (Assumption~\ref{ass:sample-pi}); a
profile paired with itself lies there with probability $\pi_j$, as it
coincides with itself by definition. The self-pairs are thus spurious
coincidences, and $n_j(n_j-1)=n_j^2-n_j$ removes exactly those $N$
terms, leaving the $N(N-1)$ ordered pairs of distinct profiles. Hence
\[
  \E\big[n_j(n_j-1)\big]=N(N-1)\,\pi_j^2 ,
  \qquad\text{so}\qquad
  \E\Big[\frac{n_j(n_j-1)}{N(N-1)}\Big]=\pi_j^2 ,
\]
the normalization matching the number of pairs retained. Each cell term
is therefore exactly unbiased for $\pi_j^2$, and \eqref{eq:bias-diag}
gives
\[
  \mathrm{Bias}\big(\sigwhatsq\big)
  =K\sum_{j=1}^{K}\big(\pi_j^2-\pi_j^2\big)=0 ,
\]
i.e. $\E\big[\sigwhatsq\big]=\sigma_{w^{\star}}^2$
for every $N\ge2$ and every $K$.
\end{proof}%
}


\begin{remark}[Trajectory dependence and the between-float estimator]
\label{rem:trajectory}
Estimator~\eqref{eq:sigmawK-hat} treats all profiles as independent draws
from $\pi$ (Assumption~\ref{ass:sample-pi}), yet profiles from one float
trace a single advective path and are not mutually independent. The
resulting error is nonetheless one-sided under the physically relevant
form of that dependence. Writing $c_a$ for the cell of profile $a$ and
extending the computation of Lemma~\ref{lem:unbiased} without its
independence step,
\begin{equation}
  \E\big[\sigwhatsq\big]-\sigma_{w^{\star}}^2
  =\frac{K}{N(N-1)}\sum_{a\ne b}
     \Big(\Prob(c_a=c_b)-\textstyle\sum_j\pi_j^2\Big),
  \label{eq:diag-bias-identity}
\end{equation}
an exact identity over the ordered pairs of distinct profiles. Two
independent draws from $\pi$ share a cell with probability exactly
$\sum_j\pi_j^2$, so each summand is the excess coincidence of a pair over
that baseline. Between-float pairs contribute nothing when distinct floats
are independent, and a non-negative amount under the shared advection
discussed below; within-float pairs---two cycles of one float, some ten
days and a short displacement apart relative to the cell scale---fall in a
common cell at least as often as independent draws, so their summands are
non-negative. Hence $\E\big[\sigwhatsq\big]\ge\sigma_{w^{\star}}^2$:
within-float clustering inflates \eqref{eq:sigmawK-hat}, the conservative
direction for an upper bound on $\sigma_{w^{\star}}$.
\end{remark}

Therefore, we also present an alternative \emph{between-float} estimator that counts only coincidence pairs from distinct floats.
Let $m_{\ell j}$ be the number of profiles of float $\ell$ in cell $j$ and $d_\ell=\sum_j m_{\ell j}$ its total, \begin{equation} \sigwhatsqbf = K\,\frac{\sum_j n_j(n_j-1)-\sum_\ell\sum_j m_{\ell j}(m_{\ell j}-1)} {N(N-1)-\sum_\ell d_\ell(d_\ell-1)}\;-\;1 , \label{eq:sigmawK-bf} \end{equation} which removes the within-float diagonal and normalizes by the between-float pair total.
Under the assumption that \emph{distinct} floats are independent draws from $\pi$, it is unbiased for $\sigma_{w^{\star}}^2$ (Lemma~\ref{lem:bf-unbiased}), with no within-float independence required, so it removes the within-float autocorrelation to which \eqref{eq:sigmawK-hat} is exposed.
This between-float independence is, however, only approximate: floats with a shared deployment or advective history---nearby releases entrained in the same current or mesoscale features---have positively associated cell memberships, so two of their profiles fall in a common cell with probability at least $\sum_j\pi_j^2$, and hence $\E\big[\sigwhatsqbf\big]\ge\sigma_{w^{\star}}^2$. The residual bias is thus in the same upward direction as the within-float clustering the estimator removes---conservative for an upper bound on $\sigma_{w^{\star}}$---and the between-float coherence it leaves is the weaker term, the same water parcel no longer being shared and trajectories diverging over time. We therefore use \eqref{eq:sigmawK-bf} in this spirit, as a conservative input to the bound rather than an exactly unbiased estimator.
Comparing \eqref{eq:sigmawK-hat} and \eqref{eq:sigmawK-bf}, which differ only through the within-float diagonal terms removed from the latter, is moreover an assumption-free diagnostic of how much of the sampling variability is trajectory autocorrelation rather than preferential sampling.

\begin{lemma}[Unbiasedness of the between-float estimator]
\label{lem:bf-unbiased}
If profiles from \emph{distinct} floats are independent draws with marginal
distribution $\pi$ then $\E\big[\sigwhatsqbf\big] = \sigma_{w^{\star}}^2$.
\end{lemma}
\deferproof{bf-unbiased}{%
\begin{proof}[Proof of Lemma~\ref{lem:bf-unbiased}]
Index profiles by $(\ell,i)$, with $\ell=1,\dots,L$ the float and
$i=1,\dots,d_\ell$ the profile within float $\ell$; let
$c_{\ell i}\in\{1,\dots,K\}$ be the cell membership of profile $(\ell,i)$, so that
\[
  m_{\ell j}=\sum_{i=1}^{d_\ell}\1\{c_{\ell i}=j\},\qquad
  n_j=\sum_{\ell}m_{\ell j},\qquad
  N=\sum_\ell d_\ell .
\]
We use only the two hypotheses of the lemma:
\begin{enumerate}
\item[(A1)] $\Prob(c_{\ell i}=j)=\pi_j$ for every profile $(\ell,i)$ and
  cell $j$;
\item[(A2)] the cell membership vectors of distinct floats are independent,
  $(c_{\ell i})_{i\le d_\ell}\perp\!\!\!\perp(c_{\ell'i'})_{i'\le d_{\ell'}}$
  for $\ell\ne\ell'$.
\end{enumerate}
Nothing is assumed about the joint distribution
$(c_{\ell 1},\dots,c_{\ell d_\ell})$ of the cells within a single float.
Write
$S_{\mathrm{bf}}=\sum_j n_j(n_j-1)-\sum_\ell\sum_j m_{\ell j}(m_{\ell j}-1)$
and $P_{\mathrm{bf}}=N(N-1)-\sum_\ell d_\ell(d_\ell-1)$ for the numerator
and denominator of~\eqref{eq:sigmawK-bf}, so
$\sigwhatsqbf=K\,S_{\mathrm{bf}}/P_{\mathrm{bf}}-1$.

The whole argument rests on one structural observation: $S_{\mathrm{bf}}$
and $P_{\mathrm{bf}}$ are sums over the \emph{same} index set, the ordered
pairs of profiles drawn from two \emph{distinct} floats. The denominator
counts all such pairs; the numerator counts those that land in a common
cell. Each pair coincides with one and the same expected probability
$\sum_j\pi_j^2$, so $\E[S_{\mathrm{bf}}]=\bigl(\sum_j\pi_j^2\bigr)
P_{\mathrm{bf}}$ term by term, and the ratio is exactly $\sum_j\pi_j^2$.
The three steps below make this precise.

\emph{Step 1 (the between-float pair set).}
Let
\[
  \mathcal B=\bigl\{\bigl((\ell,i),(\ell',i')\bigr):\ \ell\ne\ell',\
    i\le d_\ell,\ i'\le d_{\ell'}\bigr\}
\]
be the ordered pairs of profiles belonging to two distinct floats. Its
size depends only on the design $(d_\ell)_\ell$: since $N=\sum_\ell d_\ell$,
\[
  |\mathcal B|=\sum_{\ell\ne\ell'}d_\ell d_{\ell'}
  =\Bigl(\sum_\ell d_\ell\Bigr)^2-\sum_\ell d_\ell^2
  =N(N-1)-\sum_\ell d_\ell(d_\ell-1)=P_{\mathrm{bf}} .
\]
Thus $P_{\mathrm{bf}}$ is a deterministic constant, strictly positive as
soon as two floats each contribute a profile.

\emph{Step 2 ($S_{\mathrm{bf}}$ counts the coinciding pairs in $\mathcal B$).}
Expanding $n_j=\sum_\ell m_{\ell j}$ separates diagonal from off-diagonal
terms,
\[
  n_j(n_j-1)=n_j^2-n_j
  =\sum_\ell m_{\ell j}(m_{\ell j}-1)+\sum_{\ell\ne\ell'}m_{\ell j}m_{\ell'j},
\]
so summing over $j$ cancels the within-float diagonal against the term
subtracted in $S_{\mathrm{bf}}$, leaving
\[
  S_{\mathrm{bf}}=\sum_{\ell\ne\ell'}\sum_{j=1}^{K}m_{\ell j}m_{\ell'j}
  =\sum_{\ell\ne\ell'}\sum_{i=1}^{d_\ell}\sum_{i'=1}^{d_{\ell'}}
    \sum_{j=1}^{K}\1\{c_{\ell i}=j\}\1\{c_{\ell'i'}=j\}
  =\sum_{((\ell,i),(\ell',i'))\in\mathcal B}\1\{c_{\ell i}=c_{\ell'i'}\},
\]
where the second equality inserts the definition of $m_{\ell j}$ and the
third uses
$\sum_j\1\{c_{\ell i}=j\}\1\{c_{\ell'i'}=j\}=\1\{c_{\ell i}=c_{\ell'i'}\}$. So
$S_{\mathrm{bf}}$ is exactly the number of pairs in $\mathcal B$ that fall
in a common cell; every pair with $\ell=\ell'$ has dropped out.

\emph{Step 3 (each pair in $\mathcal B$ coincides with probability
$\sum_j\pi_j^2$).}
Fix $\bigl((\ell,i),(\ell',i')\bigr)\in\mathcal B$, so $\ell\ne\ell'$.
Decomposing over the common cell and applying (A2) then (A1),
\[
  \E\bigl[\1\{c_{\ell i}=c_{\ell'i'}\}\bigr]
  =\sum_{j=1}^{K}\Prob\bigl(c_{\ell i}=j,\;c_{\ell'i'}=j\bigr)
  \overset{\text{(A2)}}{=}\sum_{j=1}^{K}\Prob(c_{\ell i}=j)\Prob(c_{\ell'i'}=j)
  \overset{\text{(A1)}}{=}\sum_{j=1}^{K}\pi_j^{2}.
\]
This is the same constant for every pair in $\mathcal B$, independent of
$i,i'$; (A2) is invoked only across the two distinct floats $\ell\ne\ell'$,
never within a float.

\emph{Step 4 (assembly).}
By Step 2 and linearity, summing the constant of Step 3 over the
$|\mathcal B|=P_{\mathrm{bf}}$ pairs of Step 1,
\[
  \E[S_{\mathrm{bf}}]
  =\sum_{((\ell,i),(\ell',i'))\in\mathcal B}
     \E\bigl[\1\{c_{\ell i}=c_{\ell'i'}\}\bigr]
  =P_{\mathrm{bf}}\sum_{j=1}^{K}\pi_j^{2}.
\]
Since $P_{\mathrm{bf}}$ is deterministic and nonzero,
\[
  \E\bigl[\sigwhatsqbf\bigr]
  =\frac{K\,\E[S_{\mathrm{bf}}]}{P_{\mathrm{bf}}}-1
  =K\sum_{j=1}^{K}\pi_j^{2}-1
  =\sigma_{w^{\star}}^{2}.
\]
The within-float contributions $\sum_\ell\sum_j m_{\ell j}(m_{\ell j}-1)$ and
$\sum_\ell d_\ell(d_\ell-1)$ cancelled in Steps~1--2, before any expectation was
taken, so no assumption on the dependence among a single float's profiles
was needed. (Such dependence still affects $\Var[\sigwhatsqbf]$, but not
its mean.)
\end{proof}%
}

\subsection{A data-only upper bound on the field dispersion}
\label{sec:sigmaY-bound}

The scale-decomposed bound of Proposition~\ref{prop:scale-decomp}
requires a value for the field dispersion, over all of $\X$ and not only
where floats went. A dense external product would supply it directly,
but at a price that Remark~\ref{rem:conservative-Y} makes precise. We
bound it instead from the profiles themselves. The obstruction is that a
float-only dispersion targets the wrong measure
(Proposition~\ref{prop:float-contam}); the resolution is that this
contamination is itself a preferential-sampling bias---of the
squared-anomaly field---so the same calculus that bounds the
first-moment bias bounds it too, converting the observed OHC variance
into a conservative upper bound that requires only a single scalar shape
constant.

Everything in this subsection is stated for a generic field
$y\in L^2(\nu)$ and is applied in practice to the residual field $y'$ of
Definition~\ref{def:offset}.

\begin{proposition}[Float-only field dispersion is contaminated]
\label{prop:float-contam}
The dispersion of $Y$ over the float locations estimates a design-measure
dispersion, not $\Var_\nu(Y)=\sigma_Y^2$. Writing $g:=(Y-\mu)^2$ for the
squared anomaly about the \emph{target} mean,
\begin{equation}
  \E_\pi[g]-\Var_\nu(Y)=\Cov_\nu\!\big(w,g\big).
  \label{eq:float-contam}
\end{equation}
\end{proposition}
\deferproof{float-contam}{%
\begin{proof}[Proof of Proposition~\ref{prop:float-contam}]
The float locations are draws from $\pi$
(Assumption~\ref{ass:sample-pi}). By the change of measure
$\dd\pi=w\dd\nu$ (Definition~\ref{def:w}) and $\E_\nu[w]=1$,
\[
  \E_\pi[g]=\E_\nu[wg]=\E_\nu[g]+\E_\nu\big[(w-1)g\big]
  =\Var_\nu(Y)+\Cov_\nu(w,g),
\]
using $\E_\nu[g]=\Var_\nu(Y)$ for the middle term and
$\Cov_\nu(w,g)=\E_\nu[(w-1)(g-\E_\nu g)]=\E_\nu[(w-1)g]$ for the last.
No centring of the design measure is involved.
\end{proof}%
}

Identity \eqref{eq:float-contam} concerns the second moment about the
target mean $\mu$. The empirical variance of the observed values centres
instead at their own sample mean, and so estimates $\Var_\pi(Y)$, which
centres at the design mean $\mu_w$; the two dispersions differ by the
parallel-axis term $(\mu_w-\mu)^2=B_{\mathrm{samp}}^2$.
Proposition~\ref{prop:var-contam} records the resulting exact identity,
and it is that form which is used throughout.

The discrepancy is non-zero for exactly the reason $B_{\mathrm{samp}}$
is: it is the $w$--$Y$ coupling this paper studies, applied to
$(Y-\mu)^2$. A float-only $\hat\sigma_Y$ is thus inflated when floats
favour the high-variability front and deflated when they collect in
quiescent water. De-biasing by importance weighting
($\widehat{\sigma_Y^2}=\frac1n\sum_i w(x_i)^{-1}(y(x_i)-\hat\mu)^2$)
recovers $\Var_\nu(Y)$ in principle but requires $w$---the object we
avoid modelling---and is dominated by $1/w$ in the undersampled cells
where floats are scarcest, so its variance is largest exactly where the
bound is most sensitive. We therefore do not de-bias, but bound.

\begin{remark}[Why the field dispersion is not imported]
\label{rem:conservative-Y}
The natural alternative is to read $\sigma_Y$ off a dense reanalysis.
Data assimilation, however, tends to \emph{attenuate} anomaly amplitude,
so such a product is liable to understate $\sigma_Y$---the unsafe
direction for an upper bound, and one whose magnitude cannot be
certified from the product itself. The failure is silent: an attenuated
$\hat\sigma_Y$ returns a bound that is too small with no diagnostic
signalling it, and the natural remedy of recomputing across products
bounds only the spread between products, not their common bias. This is
the asymmetry that governs every external input in this paper. A product
may legitimately supply a dimensionless \emph{shape} (the kurtosis of
Remark~\ref{rem:kappa}, which is affine-invariant and so blind to
attenuation), a \emph{scale} (the decorrelation length, which fixes a
partition rather than a
magnitude), or the \emph{location} of a control variate
(where error costs efficiency and not
validity)---but never an \emph{amplitude} entering a bound
multiplicatively.
\end{remark}

\begin{assumption}[Finite fourth moment]
\label{ass:l4}
The realized field satisfies $y\in L^4(\nu)$, equivalently
$\sigma_g:=\big\|\,(y-\mu)^2-\sigma_Y^2\,\big\|_{L^2(\nu)}<\infty$.
\end{assumption}

This is the field-side analogue of the $w\in L^2(\nu)$ restriction of
Definition~\ref{def:w}: the first-moment bound needs two field moments,
the second-moment bound four. Throughout this section write
$g:=(y-\mu)^2$ for the \emph{squared-anomaly field}, so that
$\E_\nu[g]=\sigma_Y^2$ (Definition~\ref{def:mu}) and
$\sigma_g^2=\Var_\nu(g)$.

\begin{proposition}[Exact second-moment contamination]
\label{prop:var-contam}
Let $R:=\Var_\pi(Y)$ be the design variance of the field---the
population analogue of the empirical variance of $y$ over the float
locations. Under Assumption~\ref{ass:l4},
\[
  R=\sigma_Y^2+\Cov_\nu\!\big(w(X),\,g(X)\big)-B_{\mathrm{samp}}^2 .
\]
\end{proposition}
\deferproof{var-contam}{%
\begin{proof}[Proof of Proposition~\ref{prop:var-contam}]
Centering the design variance at the target mean $\mu$ rather than the
design mean $\mu_w$ and using $\E_\pi[y-\mu]=\mu_w-\mu=B_{\mathrm{samp}}$
(Definition~\ref{def:mu-w}, Proposition~\ref{prop:cov-identity}),
\[
  R=\E_\pi\big[(y-\mu_w)^2\big]
   =\E_\pi\big[(y-\mu)^2\big]-(\mu_w-\mu)^2
   =\E_\pi[g]-B_{\mathrm{samp}}^2 .
\]
Change of measure $\dd\pi=w\dd\nu$ (Definition~\ref{def:w}) gives
$\E_\pi[g]=\E_\nu[w\,g]$, and since $\E_\nu[w-1]=0$,
\[
  \E_\nu[w\,g]=\E_\nu[g]+\E_\nu[(w-1)g]
             =\sigma_Y^2+\Cov_\nu\!\big(w,g\big),
\]
the last equality because
$\Cov_\nu(w,g)=\E_\nu\big[(w-1)(g-\E_\nu g)\big]=\E_\nu[(w-1)g]$.
Combining the two displays yields the identity.
\end{proof}%
}

\begin{remark}[Centring, and why the squared bias is retained]
\label{rem:var-contam-order}
Propositions~\ref{prop:float-contam} and~\ref{prop:var-contam} are not
two forms of one identity, and the difference between them is not a
truncation: they concern two different dispersions. The first centres the
squared anomaly at the target mean $\mu$, the second at the design mean
$\mu_w$, and the two differ by exactly
$(\mu_w-\mu)^2=B_{\mathrm{samp}}^2$. Writing $\Var_\pi(Y)$ on the left of
\eqref{eq:float-contam} would therefore be false whenever
$B_{\mathrm{samp}}\ne0$, which is precisely the case this paper is about.

The distinction is not cosmetic, because the two forms err in opposite
directions. Since $B_{\mathrm{samp}}^2\ge0$ enters
Proposition~\ref{prop:var-contam} with a negative sign, it \emph{widens}
the envelope of Proposition~\ref{prop:sigmaY-envelope}; substituting the
$\mu$-centred identity in its place would yield
$\sigma_Y^2\le R+|\Cov_\nu(w,g)|$, a strictly smaller and hence
anti-conservative bound. The term is $O(\sigma_w^2\sigma_Y^2)$, second
order in the sampling unevenness, but it is retained because it is
required for validity and not merely for accuracy.
\end{remark}

Because $\Cov_\nu(w,g)$ is the covariance of the weight with the field
$g$, it is a preferential-sampling bias in the exact sense of
Proposition~\ref{prop:cov-identity}, now carried by the squared-anomaly
field. Every step of Sections~\ref{sec:defbias}--\ref{sec:estimation}
therefore transfers verbatim with $g$ in place of $y$. Decomposing
$g-\sigma_Y^2=g^{\star}+g^{\perp}$ through the split~\eqref{eq:split}
and repeating Proposition~\ref{prop:scale-decomp},
\begin{equation}
  \big|\Cov_\nu(w,g)\big|
  \;\le\;\sigma_{w^{\star}}\sigma_{g^{\star}}
        +\sigma_{w^{\perp}}\sigma_{g^{\perp}}
  \;\le\;\sigma_w\sigma_g ,
  \label{eq:var-scale-decomp}
\end{equation}
where
$\sigma_{g^{\star}}:=\|\PK(g-\sigma_Y^2)\|_{L^2(\nu)}$ and
$\sigma_{g^{\perp}}:=\|\PKp(g-\sigma_Y^2)\|_{L^2(\nu)}$.

\begin{lemma}[Kurtosis reduction]
\label{lem:kurt}
Let $\kappa:=\E_\nu[(y-\mu)^4]/\sigma_Y^4$ be the reference-measure
kurtosis of the anomaly field. Then $\kappa\ge1$ and
\[
  \sigma_g=\sigma_Y^2\sqrt{\kappa-1}.
\]
\end{lemma}
\deferproof{kurt}{%
\begin{proof}[Proof of Lemma~\ref{lem:kurt}]
By definition of $g$ and $\E_\nu[g]=\sigma_Y^2$,
\[
  \sigma_g^2=\Var_\nu(g)=\E_\nu[g^2]-\big(\E_\nu g\big)^2
   =\E_\nu\big[(y-\mu)^4\big]-\sigma_Y^4=(\kappa-1)\,\sigma_Y^4 .
\]
Non-negativity of the variance gives $\kappa\ge1$ (equivalently, Jensen's
inequality applied to $G=(y-\mu)^2$).
\end{proof}%
}

The one factor of \eqref{eq:var-scale-decomp} not paired with an
estimable weight dispersion is the within-cell term
$\sigma_{w^{\perp}}\sigma_{g^{\perp}}$, and it is controlled by exactly
the cell-scale choice that already controls the mean-bound residual,
because a field that is nearly constant on a cell has a nearly constant
square.

\begin{lemma}[Projection contracts fourth moments]
\label{lem:l4-contraction}
$\PK$ is the conditional expectation given the $\sigma$-algebra generated
by $\mathcal{P}_K$, hence a contraction on $L^p(\nu)$ for every
$p\ge1$. In particular
\[
  \E_\nu\big[(y^{\star})^4\big]\;\le\;\E_\nu\big[(y-\mu)^4\big]
  =\kappa\,\sigma_Y^4,
  \qquad\text{equivalently}\qquad
  \big(\kappa^{\star}\big)^{1/4}\sigma_{Y^{\star}}
  \;\le\;\kappa^{1/4}\sigma_Y ,
\]
where $\kappa^{\star}:=\E_\nu[(y^{\star})^4]/\sigma_{Y^{\star}}^4$.
\end{lemma}
\deferproof{l4-contraction}{%
\begin{proof}[Proof of Lemma~\ref{lem:l4-contraction}]
Cells have equal reference measure, so $\PK f$ is the $\nu$-conditional
expectation of $f$ given cell membership. Conditional Jensen applied to
the convex map $z\mapsto z^4$ gives
$(\E_\nu[f\mid\mathcal{P}_K])^4\le\E_\nu[f^4\mid\mathcal{P}_K]$
pointwise; taking $\E_\nu$ of both sides and $f=y-\mu$ yields the first
inequality. The second is the first rewritten using
$\E_\nu[(y^{\star})^4]=\kappa^{\star}\sigma_{Y^{\star}}^4$ and taking
fourth roots.
\end{proof}%
}

The practical value of Lemma~\ref{lem:l4-contraction} is that it removes
a declared constant: the between-cell kurtosis never has to be specified,
because it appears only in the combination the lemma bounds.

\begin{lemma}[The within-cell residual is governed by within-cell field
variation]
\label{lem:gperp}
Let $\kappa^{\perp}:=\E_\nu[(y^{\perp})^4]/\sigma_{Y^{\perp}}^4$ be the
kurtosis of the within-cell component. Then
\begin{equation}
  \sigma_{g^{\perp}}\;\le\;c_g\,\sigma_{Y^{\perp}},
  \qquad
  c_g:=2\,\sigma_Y\big(\kappa\,\kappa^{\perp}\big)^{1/4}
       +\sigma_{Y^{\perp}}\sqrt{\kappa^{\perp}} .
  \label{eq:gperp-bound}
\end{equation}
In particular $\sigma_{g^{\perp}}\to0$ as $\sigma_{Y^{\perp}}\to0$ under
refinement, so the within-cell residual of \eqref{eq:var-scale-decomp} is
controlled by the same cell-scale choice that controls the mean-bound
residual $\sigma_{w^{\perp}}\sigma_{Y^{\perp}}$.
\end{lemma}
\deferproof{gperp}{%
\begin{proof}[Proof of Lemma~\ref{lem:gperp}]
On a cell $C_j$ set $a_j:=\bar y_j-\mu$ and $y^{\perp}:=y-\bar y_j$, so
$y-\mu=a_j+y^{\perp}$ and $g=a_j^2+2a_jy^{\perp}+(y^{\perp})^2$. Since
$\int_{C_j}y^{\perp}\dd\nu=0$, the cell mean is
$\bar g_j=a_j^2+\overline{(y^{\perp})^2}_j$, whence
\[
  g^{\perp}=g-\bar g_j
   =2a_jy^{\perp}+\Big((y^{\perp})^2-\overline{(y^{\perp})^2}_j\Big).
\]
Taking $L^2(\nu)$ norms and applying the triangle inequality,
$\sigma_{g^{\perp}}\le 2\|y^{\star}y^{\perp}\|_{L^2(\nu)}
 +\|\PKp((y^{\perp})^2)\|_{L^2(\nu)}$, where $y^{\star}$ denotes the
piecewise-constant function equal to $a_j$ on $C_j$. For the first term,
Cauchy--Schwarz in $L^2(\nu)$ applied to $(y^{\star})^2$ and
$(y^{\perp})^2$ gives
\[
  \|y^{\star}y^{\perp}\|_{L^2(\nu)}^2
  =\E_\nu\big[(y^{\star})^2(y^{\perp})^2\big]
  \le\Big(\E_\nu[(y^{\star})^4]\,\E_\nu[(y^{\perp})^4]\Big)^{1/2}
  \le\big(\kappa\,\kappa^{\perp}\big)^{1/2}
   \sigma_Y^2\sigma_{Y^{\perp}}^2 ,
\]
the last step by Lemma~\ref{lem:l4-contraction} together with
$\E_\nu[(y^{\perp})^4]=\kappa^{\perp}\sigma_{Y^{\perp}}^4$.
For the second, $\PKp$ is a contraction (Lemma~\ref{lem:projection}), so
$\|\PKp((y^{\perp})^2)\|\le\|(y^{\perp})^2\|_{L^2(\nu)}
 =\sqrt{\kappa^{\perp}}\,\sigma_{Y^{\perp}}^2$. Adding the two gives
\eqref{eq:gperp-bound}.
\end{proof}%
}

\begin{remark}[Why moments rather than a supremum]
\label{rem:gperp-moments}
An earlier form of \eqref{eq:gperp-bound} used $\|y^{\star}\|_\infty$ in
place of the moment factor $\kappa^{1/4}\sigma_Y$. The moment form is
preferred because a supremum cannot be estimated safely from a finite
sample: the maximum over $F$ floats systematically \emph{under}-estimates
$\|y^{\star}\|_\infty$, which is the unsafe direction for an upper bound,
whereas fourth moments are estimated without such one-sided error.
\end{remark}

We now carry the within-cell residuals through the argument rather than
discarding them. Both enter only through the single scalar
\begin{equation}
  \varepsilon\;:=\;\sigma_{w^{\perp}}\,\sigma_{Y^{\perp}},
  \label{eq:residual-budget}
\end{equation}
the \emph{residual budget}, which is exactly the quantity already
appearing in the mean bound of Proposition~\ref{prop:scale-decomp}. No
new unknown is introduced.

\begin{proposition}[Data-only upper bound on $\sigma_Y$]
\label{prop:sigmaY-envelope}
Set $\Delta:=\sigma_{w^{\star}}\sqrt{\kappa-1}+\sigma_{w^{\star}}^2$ and
suppose $\Delta<1$. Then for the residual budget $\varepsilon$ of
\eqref{eq:residual-budget},
\begin{equation}
  \sigma_Y\;\le\;\overline{\sigma}_Y(\varepsilon)
  \;:=\;\frac{b\,\varepsilon
    +\sqrt{\,b^2\varepsilon^2
      +(1-\Delta)\big(R+d\,\varepsilon+\varepsilon^2\big)\,}}
    {1-\Delta},
  \label{eq:envelope-eps}
\end{equation}
where $b:=\sigma_{w^{\star}}+(\kappa\,\kappa^{\perp})^{1/4}$ and
$d:=\sigma_{Y^{\perp}}\sqrt{\kappa^{\perp}}$. Only two shape constants
enter: $\kappa$, which is declared, and $\kappa^{\perp}$, which is
estimated from within-cell float pairs by
\eqref{eq:kappaperp-pairs}. The map
$\varepsilon\mapsto\overline{\sigma}_Y(\varepsilon)$ is continuous and
strictly increasing, and
\[
  \overline{\sigma}_Y(0)=\sqrt{\frac{R}{1-\Delta}} .
\]
\end{proposition}
\deferproof{sigmaY-envelope}{%
\begin{proof}[Proof of Proposition~\ref{prop:sigmaY-envelope}]
Proposition~\ref{prop:var-contam} gives exactly
$\sigma_Y^2=R-\Cov_\nu(w,g)+B_{\mathrm{samp}}^2
 \le R+|\Cov_\nu(w,g)|+B_{\mathrm{samp}}^2$. For the covariance, apply
\eqref{eq:var-scale-decomp}, then $\sigma_{g^{\star}}\le\sigma_g$
(Lemma~\ref{lem:coarsening} applied to $g$) with
Lemma~\ref{lem:kurt}, and Lemma~\ref{lem:gperp} for the within-cell part:
\[
  |\Cov_\nu(w,g)|
  \le\sigma_{w^{\star}}\sqrt{\kappa-1}\,\sigma_Y^2
    +c_g\,\sigma_{w^{\perp}}\sigma_{Y^{\perp}}
  =\sigma_{w^{\star}}\sqrt{\kappa-1}\,\sigma_Y^2+c_g\varepsilon .
\]
By \eqref{eq:gperp-bound},
$c_g\varepsilon=2(\kappa\,\kappa^{\perp})^{1/4}\varepsilon\,\sigma_Y
 +d\varepsilon$ exactly, with no further substitution. For the squared
bias,
Proposition~\ref{prop:scale-decomp} and the same inequality give
\[
  B_{\mathrm{samp}}^2
  \le\big(\sigma_{w^{\star}}\sigma_{Y^{\star}}+\varepsilon\big)^2
  \le\sigma_{w^{\star}}^2\sigma_Y^2
     +2\sigma_{w^{\star}}\varepsilon\,\sigma_Y+\varepsilon^2 .
\]
Summing and collecting powers of $\sigma_Y$,
\[
  (1-\Delta)\,\sigma_Y^2-2b\varepsilon\,\sigma_Y
  -\big(R+d\varepsilon+\varepsilon^2\big)\;\le\;0 .
\]
Under $\Delta<1$ the leading coefficient is positive, so the admissible
set is the interval between the roots and $\sigma_Y$ is at most the
larger one, which is \eqref{eq:envelope-eps}. Monotonicity and continuity
in $\varepsilon$ are immediate from the formula, and setting
$\varepsilon=0$ gives $\sqrt{(1-\Delta)R}/(1-\Delta)$.
\end{proof}%
}

Every quantity in \eqref{eq:envelope-eps} except $\varepsilon$ is
available without a field product: $R$ from the observed OHC,
$\sigma_{w^{\star}}$ from the count statistic
$\sigwhat$~\eqref{eq:sigmawK-hat}, $\sigma_{Y^{\perp}}$ from within-cell
float pairs, and
$\kappa$ a chosen constant and $\kappa^{\perp}$ estimated
(Remark~\ref{rem:kappa}). The residual budget $\varepsilon$ is not
estimable---that is the standing conclusion of
Section~\ref{sec:decomp-wy}---and so is never assumed small. It is
carried as an argument, and the bound is reported as a function of it.

\begin{remark}[When $\Delta\ge1$]
\label{rem:delta-fails}
The hypothesis $\Delta<1$ is not a technical convenience and cannot be
repaired. When $\Delta\ge1$ the inequality
$\sigma_Y^2\le R+\Delta\sigma_Y^2+\cdots$ is satisfied by arbitrarily
large $\sigma_Y$, and the failure is informational rather than
algebraic: at that level of sampling unevenness the second-moment
contamination could account for the entire observed dispersion, so the
spread of the float values is consistent with a quiet field sampled very
unevenly and with a variable field sampled evenly. No estimator can
separate the two from the profiles alone, and the only remedy is to
coarsen the partition until $\sigma_{w^{\star}}$ falls below the ceiling
$\sigma_{\max}(\kappa)$ at which the envelope \eqref{eq:envelope-eps}
remains finite. Because $\sigwhat$ is itself
estimated, the check should use the upper limit of its jackknife interval
rather than the point estimate, and the same $\kappa$ must be used in the
ceiling and in \eqref{eq:envelope-eps}: a large $\kappa$ chosen for
conservatism in the field factor cannot be paired with a small one chosen
to license a finer partition.
\end{remark}

\begin{corollary}[Field-product-free bias bound]
\label{cor:field-free}
Under the hypotheses of Proposition~\ref{prop:sigmaY-envelope}, for every
residual budget $\varepsilon$,
\begin{equation}
  |B_{\mathrm{samp}}|
  \;\le\;\mathcal{B}(\varepsilon)
  \;:=\;\sigma_{w^{\star}}\,\overline{\sigma}_Y(\varepsilon)
        +\varepsilon ,
  \label{eq:bias-eps}
\end{equation}
and $\mathcal{B}$ is continuous and strictly increasing in $\varepsilon$.
\end{corollary}
\deferproof{field-free}{%
\begin{proof}[Proof of Corollary~\ref{cor:field-free}]
Proposition~\ref{prop:scale-decomp} gives
$|B_{\mathrm{samp}}|\le\sigma_{w^{\star}}\sigma_{Y^{\star}}
 +\sigma_{w^{\perp}}\sigma_{Y^{\perp}}$. Replace $\sigma_{Y^{\star}}$ by
the larger $\sigma_Y\le\overline{\sigma}_Y(\varepsilon)$
(Lemma~\ref{lem:coarsening} and
Proposition~\ref{prop:sigmaY-envelope}) and the second term by
$\varepsilon$ by definition \eqref{eq:residual-budget}. Monotonicity is
inherited from $\overline{\sigma}_Y$.
\end{proof}%
}

\subsubsection*{Negligibility, made quantitative}
\label{sec:breakdown}

Since $\mathcal{B}$ is increasing and continuous, the claim ``the
residual is negligible'' can be replaced by a number.

\begin{definition}[Breakdown budget]
\label{def:breakdown}
Fix a benchmark $\mathcal{T}>0$ in the units of the estimand. The
\emph{breakdown budget} is
\[
  \varepsilon^{\star}(\mathcal{T})
  :=\sup\{\varepsilon\ge0:\ \mathcal{B}(\varepsilon)\le\mathcal{T}\},
\]
with $\varepsilon^{\star}=0$ if $\mathcal{B}(0)>\mathcal{T}$. The
reported conclusion is then not that $\varepsilon$ is small but that the
bound remains below $\mathcal{T}$ for every residual budget up to
$\varepsilon^{\star}(\mathcal{T})$.
\end{definition}

This converts an adjective into a falsifiable statement: a reader who
believes the sub-cell weight variation could exceed
$\varepsilon^{\star}/\sigma_{Y^{\perp}}$ has a concrete objection, and
one that can be argued from the float geometry. We report
$\mathcal{B}(\varepsilon)$ as a curve, together with
$\varepsilon^{\star}$ at two benchmarks: the OHC trend over the analysis
window, and the posterior standard deviation of the OHC estimate the
bound is meant to qualify. Both are in the same units as
$B_{\mathrm{samp}}$ and are the quantities against which a bias is
decision-relevant.

\begin{remark}[Closing the budget by declaring a total unevenness]
\label{rem:chi2-ceiling}
If a single number is wanted in place of a curve, $\varepsilon$ can be
eliminated at the cost of one explicit assumption. Since
$\sigma_{w^{\perp}}^2=\sigma_w^2-\sigma_{w^{\star}}^2$
(Lemma~\ref{lem:coarsening}), a declared ceiling
$\sigma_w^2=\chi^2(\pi\|\nu)\le\chi^2_{\max}$ on the \emph{total}
sampling unevenness gives
\[
  \varepsilon\;\le\;\sqrt{\big(\chi^2_{\max}-\sigma_{w^{\star}}^2\big)_+}
  \;\cdot\;\sigYperphat ,
\]
and \eqref{eq:bias-eps} evaluated there is a closed bound. We prefer the
curve as the primary report, since $\chi^2_{\max}$ is an assumption the
framework otherwise avoids; but it is an interpretable and criticizable
one---a statement about how non-uniform Argo sampling could possibly
be---and is stated as a variant for readers who want a single figure.
\end{remark}

\begin{remark}[Interpreting and choosing the shape constants]
\label{rem:kappa}
Two dimensionless shape constants enter the envelope: $\kappa$, through
Lemma~\ref{lem:kurt}, and $\kappa^{\perp}$, through
Lemma~\ref{lem:gperp}. Only the first is declared; the second is
estimated from the profiles. A third, the between-cell kurtosis
$\kappa^{\star}$, is eliminated by Lemma~\ref{lem:l4-contraction} and
never has to be chosen.

\paragraph*{The leading constant.}
The constant $\kappa$ is the kurtosis of the \emph{one-point value
distribution} of the anomaly field: pool the values $y(x)-\mu$ over all
$x\in\X$ under the reference measure $\nu$ and take the standardized
fourth moment. It is not a temporal kurtosis of a series, nor a kurtosis
across ocean realizations, and by Lemma~\ref{lem:kurt} it is
forced---no summary of the field other than $\kappa$ enters the variance
term. Physically it measures \emph{intermittency}: how concentrated the
anomaly amplitude is across space--time. Rings and the meandering jet
core are compact, high-amplitude features on a quiescent background,
which fatten the tails and lift $\kappa$ above the Gaussian value $3$.

Two points guide a conservative choice. First, $\kappa$ is \emph{not}
monotone in ``frontal-ness.'' A front expressed as rare sharp excursions
on a background is leptokurtic ($\kappa>3$); a front that partitions the
box into two bulk water masses gives a \emph{bimodal}, platykurtic
distribution ($\kappa<3$, with floor $\kappa=1$). Which regime holds is
set by the de-trending choice: removing the
standing planar gradient strips the two-mass partition and raises
$\kappa$; leaving it lowers $\kappa$. So $\kappa$ must be read from the
same anomaly field that enters the bound. Reference anchors are
exact---Gaussian $3$, Laplace $6$, exponential $9$---and vertically
integrated OHC tends to the milder leptokurtic end, the depth integral
smoothing the intermittency of instantaneous surface fields.
Conservative here means \emph{large}: a larger $\kappa$ shrinks the
factor $1-\Delta$ and widens the envelope, the only ceiling being
finiteness of \eqref{eq:envelope-eps}, namely
\[
  \kappa\;<\;1+\Big(\frac{1-\sigma_{w^{\star}}^2}{\sigma_{w^{\star}}}\Big)^{2},
\]
the same condition as $\sigma_{w^{\star}}<\sigma_{\max}(\kappa)$ read in
the other direction. It lies far above any physical value when
$\sigma_{w^{\star}}$ is small, but it is not vacuous at fine partitions
(Remark~\ref{rem:delta-fails}).

Second, $\kappa$ is a far safer quantity to import than $\sigma_Y$
itself, for a structural reason: it is affine-invariant. Under any
rescaling $y\mapsto ay+b$ both $\E_\nu[(y-\mu)^4]$ and $\sigma_Y^4$ carry
a factor $a^4$ and cancel. The specific defect that made the reanalysis
unsafe as a source of $\sigma_Y$---assimilation \emph{attenuates}
anomaly amplitude (Remark~\ref{rem:conservative-Y}), the unsafe
direction---is a roughly multiplicative rescaling, to which $\kappa$ is
blind. A reanalysis, or a heavier-tailed physical prior, is therefore a
defensible source for the \emph{shape} constant $\kappa$ even though it
is not for the \emph{amplitude} $\sigma_Y$; the one caveat is that
assimilation may damp extremes more than the bulk, understating $\kappa$,
so a margin is added rather than taking a reanalysis value at face value.

\paragraph*{The residual constant.}
The constant $\kappa^{\perp}$ is the kurtosis of the within-cell
component $y^{\perp}=\PKp(y-\mu)$, evaluated on the operating partition.
Its role is strictly subordinate. It enters \eqref{eq:envelope-eps} only
through $b=\sigma_{w^{\star}}+(\kappa\,\kappa^{\perp})^{1/4}$ and
$d=\sigma_{Y^{\perp}}\sqrt{\kappa^{\perp}}$, both of which multiply the
residual budget $\varepsilon$; at $\varepsilon=0$ it leaves the bound
entirely. In particular it does \emph{not} enter $\Delta$, so it does not
affect the refinement ceiling, and
the resolution trade-off is governed
by $\kappa$ alone. It also enters weakly: the fourth root means that
doubling it inflates $(\kappa\,\kappa^{\perp})^{1/4}$ by only
$2^{1/4}\approx1.19$, and less than that in $b$, which also carries
$\sigma_{w^{\star}}$.

No ordering with $\kappa$ may be assumed. Kurtosis is not monotone under
orthogonal projection, so $\kappa^{\perp}\le\kappa$ fails in general, and
the two constants must be fixed independently rather than deduced from
one another. (The corresponding between-cell quantity is exempt: it is
not a free constant but is dominated through
Lemma~\ref{lem:l4-contraction}, which is available because $\PK$ is a
conditional expectation and $\PKp$ is not.) The constant $\kappa^{\perp}$
is
resolution-dependent at both ends: in the coarse limit
$y^{\perp}\to y-\mu$ and $\kappa^{\perp}\to\kappa$, while below the
decorrelation length $y^{\perp}$ is approximately a local gradient times
a displacement, which within any single cell is \emph{platykurtic}.
Pooling across cells, however, makes $y^{\perp}$ a scale mixture whose
mixing variable is the local gradient magnitude, and scale mixtures
inflate kurtosis; it is thus the intermittency of the \emph{gradient}
field---sharp at the front, weak in the interior---rather than of the
field itself that sets $\kappa^{\perp}$.

Unlike $\kappa$, the constant $\kappa^{\perp}$ can be estimated directly
from the profiles, using the same within-cell float pairs that supply
$\sigYperphat$. Writing $D_{ii'}:=y'(Z_i)-y'(Z_{i'})$ for pairs falling
in a common cell of the operating partition,
\begin{equation}
  \hat\kappa^{\perp}
  \;=\;\frac{2\,\overline{D^4}}{\big(\overline{D^2}\big)^{2}}\;-\;3,
  \qquad
  \overline{D^{2}}\big/2=\sigYperphatsq ,
  \label{eq:kappaperp-pairs}
\end{equation}
where the bars denote averages over within-cell pairs. The identity
underlying \eqref{eq:kappaperp-pairs} is that for two independent draws
from a common law, $\E[D^4]=2\E[y^{\perp4}]+6\sigma^4$ and
$\E[D^2]=2\sigma^2$; the second equality in \eqref{eq:kappaperp-pairs}
shows it reuses exactly the within-cell pair statistic
$\overline{D^2}/2=\sigYperphatsq$, at no additional data cost. Like every float-based quantity it targets the
design law rather than $\nu$, but it is a dimensionless shape and so
inherits the affine-invariance argument above; and since only an upper
bound is required, the estimate is padded rather than taken at face
value. Its variance is high---fourth moments of pair differences on a
thin sample---so it is used to check the plausibility of a chosen value
rather than to supply one.

Given the insensitivity noted above, a single declared value may be used
for both, setting $\kappa^{\perp}=\kappa$ and reporting
$\hat\kappa^{\perp}$ from \eqref{eq:kappaperp-pairs} only as a check.
This is the configuration we adopt: one constant is chosen, everything
else in \eqref{eq:envelope-eps} is measured. The sensitivity analysis
nevertheless varies $\kappa$ and $\kappa^{\perp}$ separately, since only
the former interacts with the choice of resolution.
\end{remark}

\begin{remark}[A one-constant variant]
\label{rem:one-constant}
The within-cell kurtosis can be dispensed with altogether, at a price
worth stating. Since $\PKp$ is a contraction
(Lemma~\ref{lem:projection}), $\sigma_{g^{\perp}}\le\sigma_g$ directly,
without Lemma~\ref{lem:gperp}. Writing $\delta:=\sigma_{w^{\perp}}$ and
$\varepsilon=\delta\,\sigma_{Y^{\perp}}$ as before,
\eqref{eq:var-scale-decomp} and Lemma~\ref{lem:kurt} give
$|\Cov_\nu(w,g)|\le(\sigma_{w^{\star}}+\delta)\sqrt{\kappa-1}\,\sigma_Y^2$,
and the argument of Proposition~\ref{prop:sigmaY-envelope} goes through
unchanged with $b$ replaced by $\sigma_{w^{\star}}$, $d$ by $0$, and
\[
  \Delta\;\longmapsto\;
  \Delta_\delta:=(\sigma_{w^{\star}}+\delta)\sqrt{\kappa-1}
    +\sigma_{w^{\star}}^2 .
\]
Only $\kappa$ appears. The cost is structural rather than numerical: the
residual has moved from the numerator of \eqref{eq:envelope-eps}, where
it was damped by the estimable ratio
$\sigma_{Y^{\perp}}/\sigma_Y$, into the denominator, where it is
undamped and erodes the refinement ceiling: since
$\Delta_\delta\le\Delta(\sigma_{w^{\star}}+\delta)$, the condition
$\sigma_{w^{\star}}<\sigma_{\max}(\kappa)-\delta$ suffices, so the
admissible band narrows by at least $\delta$. More seriously, the two versions lean on
different mechanisms for the residual to be small: the two-constant bound
requires $\sigma_{Y^{\perp}}\to0$ under refinement, which the within-cell
pair estimate $\sigYperphat$ verifies, whereas the one-constant bound
requires $\sigma_{w^{\perp}}\to0$, which by the standing argument of
Section~\ref{sec:decomp-wy} cannot be checked. Trading a measurable
damping factor for an unmeasurable one is unattractive in a bound whose
premise is that $w$ is not modelled, so we report the one-constant form
as a cross-check rather than as the primary result.
\end{remark}

\paragraph{Estimation and uncertainty.} The design variance is estimated
by the ordinary empirical variance of the observed OHC,
\[
  \hat R=\frac1N\sum_{i=1}^{N}\big(y(x_i)-\hat\mu_\pi\big)^2 ,
  \qquad
  \hat\mu_\pi=\frac1N\sum_{i=1}^{N}y(x_i),
\]
which is consistent for $\Var_\pi(Y)$ under
Assumption~\ref{ass:sample-pi} ($\hat\mu_\pi$ estimates the design mean
$\mu_w$). Its float-clustered uncertainty is propagated by the same
delete-one-float jackknife used for $\sigwhat$, and the envelope
$\overline{\sigma}_Y$ inherits its interval by monotone propagation of
the jackknife bounds on $\hat R$ and $\sigwhat$ through
Proposition~\ref{prop:sigmaY-envelope} at the chosen $\kappa$. In the
reported bound at each resolution is
$\sigwhat\,\overline{\sigma}_Y$. Where an auxiliary field product is
available it may be used to sanity-check $\overline{\sigma}_Y$ and,
through Remark~\ref{rem:kappa}, the chosen $\kappa$---but as a
cross-check read in the knowledge that it attenuates, never as an input
to the guarantee.

